\documentclass[a4paper,12pt,twoside]{memoir}

% Castellano
\usepackage[spanish,es-tabla]{babel}
\selectlanguage{spanish}
\usepackage[utf8]{inputenc}
\usepackage[T1]{fontenc}
\usepackage{lmodern} % scalable font
\usepackage{microtype}
\usepackage{placeins}

\RequirePackage{booktabs}
\RequirePackage[table]{xcolor}
\RequirePackage{xtab}
\RequirePackage{multirow}

% Links
\PassOptionsToPackage{hyphens}{url}\usepackage[colorlinks]{hyperref}
\hypersetup{
	allcolors = {red}
}

% Ecuaciones
\usepackage{amsmath}

% Rutas de fichero / paquete
\newcommand{\ruta}[1]{{\sffamily #1}}

% Párrafos
\nonzeroparskip

% Huérfanas y viudas
\widowpenalty100000
\clubpenalty100000

% Evitar solapes en el header
\nouppercaseheads

% Imagenes
\usepackage{graphicx}
\newcommand{\imagen}[2]{
	\begin{figure}[!h]
		\centering
		\includegraphics[width=0.9\textwidth]{#1}
		\caption{#2}\label{fig:#1}
	\end{figure}
	\FloatBarrier
}

\newcommand{\imagenflotante}[2]{
	\begin{figure}%[!h]
		\centering
		\includegraphics[width=0.9\textwidth]{#1}
		\caption{#2}\label{fig:#1}
	\end{figure}
}



% El comando \figura nos permite insertar figuras comodamente, y utilizando
% siempre el mismo formato. Los parametros son:
% 1 -> Porcentaje del ancho de página que ocupará la figura (de 0 a 1)
% 2 --> Fichero de la imagen
% 3 --> Texto a pie de imagen
% 4 --> Etiqueta (label) para referencias
% 5 --> Opciones que queramos pasarle al \includegraphics
% 6 --> Opciones de posicionamiento a pasarle a \begin{figure}
\newcommand{\figuraConPosicion}[6]{%
  \setlength{\anchoFloat}{#1\textwidth}%
  \addtolength{\anchoFloat}{-4\fboxsep}%
  \setlength{\anchoFigura}{\anchoFloat}%
  \begin{figure}[#6]
    \begin{center}%
      \Ovalbox{%
        \begin{minipage}{\anchoFloat}%
          \begin{center}%
            \includegraphics[width=\anchoFigura,#5]{#2}%
            \caption{#3}%
            \label{#4}%
          \end{center}%
        \end{minipage}
      }%
    \end{center}%
  \end{figure}%
}

%
% Comando para incluir imágenes en formato apaisado (sin marco).
\newcommand{\figuraApaisadaSinMarco}[5]{%
  \begin{figure}%
    \begin{center}%
    \includegraphics[angle=90,height=#1\textheight,#5]{#2}%
    \caption{#3}%
    \label{#4}%
    \end{center}%
  \end{figure}%
}
% Para las tablas
\newcommand{\otoprule}{\midrule [\heavyrulewidth]}
%
% Nuevo comando para tablas pequeñas (menos de una página).
\newcommand{\tablaSmall}[5]{%
 \begin{table}
  \begin{center}
   \rowcolors {2}{gray!35}{}
   \begin{tabular}{#2}
    \toprule
    #4
    \otoprule
    #5
    \bottomrule
   \end{tabular}
   \caption{#1}
   \label{tabla:#3}
  \end{center}
 \end{table}
}

%
%Para el float H de tablaSmallSinColores
\usepackage{float}

%
% Nuevo comando para tablas pequeñas (menos de una página).
\newcommand{\tablaSmallSinColores}[5]{%
 \begin{table}[H]
  \begin{center}
   \begin{tabular}{#2}
    \toprule
    #4
    \otoprule
    #5
    \bottomrule
   \end{tabular}
   \caption{#1}
   \label{tabla:#3}
  \end{center}
 \end{table}
}

\newcommand{\tablaApaisadaSmall}[5]{%
\begin{landscape}
  \begin{table}
   \begin{center}
    \rowcolors {2}{gray!35}{}
    \begin{tabular}{#2}
     \toprule
     #4
     \otoprule
     #5
     \bottomrule
    \end{tabular}
    \caption{#1}
    \label{tabla:#3}
   \end{center}
  \end{table}
\end{landscape}
}

%
% Nuevo comando para tablas grandes con cabecera y filas alternas coloreadas en gris.
\newcommand{\tabla}[6]{%
  \begin{center}
    \tablefirsthead{
      \toprule
      #5
      \otoprule
    }
    \tablehead{
      \multicolumn{#3}{l}{\small\sl continúa desde la página anterior}\\
      \toprule
      #5
      \otoprule
    }
    \tabletail{
      \hline
      \multicolumn{#3}{r}{\small\sl continúa en la página siguiente}\\
    }
    \tablelasttail{
      \hline
    }
    \bottomcaption{#1}
    \rowcolors {2}{gray!35}{}
    \begin{xtabular}{#2}
      #6
      \bottomrule
    \end{xtabular}
    \label{tabla:#4}
  \end{center}
}

%
% Nuevo comando para tablas grandes con cabecera.
\newcommand{\tablaSinColores}[6]{%
  \begin{center}
    \tablefirsthead{
      \toprule
      #5
      \otoprule
    }
    \tablehead{
      \multicolumn{#3}{l}{\small\sl continúa desde la página anterior}\\
      \toprule
      #5
      \otoprule
    }
    \tabletail{
      \hline
      \multicolumn{#3}{r}{\small\sl continúa en la página siguiente}\\
    }
    \tablelasttail{
      \hline
    }
    \bottomcaption{#1}
    \begin{xtabular}{#2}
      #6
      \bottomrule
    \end{xtabular}
    \label{tabla:#4}
  \end{center}
}

%
% Nuevo comando para tablas grandes sin cabecera.
\newcommand{\tablaSinCabecera}[5]{%
  \begin{center}
    \tablefirsthead{
      \toprule
    }
    \tablehead{
      \multicolumn{#3}{l}{\small\sl continúa desde la página anterior}\\
      \hline
    }
    \tabletail{
      \hline
      \multicolumn{#3}{r}{\small\sl continúa en la página siguiente}\\
    }
    \tablelasttail{
      \hline
    }
    \bottomcaption{#1}
  \begin{xtabular}{#2}
    #5
   \bottomrule
  \end{xtabular}
  \label{tabla:#4}
  \end{center}
}



\definecolor{cgoLight}{HTML}{EEEEEE}
\definecolor{cgoExtralight}{HTML}{FFFFFF}

%
% Nuevo comando para tablas grandes sin cabecera.
\newcommand{\tablaSinCabeceraConBandas}[5]{%
  \begin{center}
    \tablefirsthead{
      \toprule
    }
    \tablehead{
      \multicolumn{#3}{l}{\small\sl continúa desde la página anterior}\\
      \hline
    }
    \tabletail{
      \hline
      \multicolumn{#3}{r}{\small\sl continúa en la página siguiente}\\
    }
    \tablelasttail{
      \hline
    }
    \bottomcaption{#1}
    \rowcolors[]{1}{cgoExtralight}{cgoLight}

  \begin{xtabular}{#2}
    #5
   \bottomrule
  \end{xtabular}
  \label{tabla:#4}
  \end{center}
}




\graphicspath{ {./img/} }

% Capítulos
\chapterstyle{bianchi}
\newcommand{\capitulo}[2]{
	\setcounter{chapter}{#1}
	\setcounter{section}{0}
	\setcounter{figure}{0}
	\setcounter{table}{0}
	\chapter*{#2}
	\addcontentsline{toc}{chapter}{#2}
	\markboth{#2}{#2}
}

% Apéndices
\renewcommand{\appendixname}{Apéndice}
\renewcommand*\cftappendixname{\appendixname}

\newcommand{\apendice}[1]{
	%\renewcommand{\thechapter}{A}
	\chapter{#1}
}

\renewcommand*\cftappendixname{\appendixname\ }

% Formato de portada
\makeatletter
\usepackage{xcolor}
\newcommand{\tutor}[1]{\def\@tutor{#1}}
\newcommand{\course}[1]{\def\@course{#1}}
\definecolor{cpardoBox}{HTML}{E6E6FF}
\def\maketitle{
  \null
  \thispagestyle{empty}
  % Cabecera ----------------
\noindent\includegraphics[width=\textwidth]{cabecera}\vspace{1cm}%
  \vfill
  % Título proyecto y escudo informática ----------------
  \colorbox{cpardoBox}{%
    \begin{minipage}{.8\textwidth}
      \vspace{.5cm}\Large
      \begin{center}
      \textbf{TFG del Grado en Ingeniería Informática}\vspace{.6cm}\\
      \textbf{\LARGE\@title{}}
      \end{center}
      \vspace{.2cm}
    \end{minipage}

  }%
  \hfill\begin{minipage}{.20\textwidth}
    \includegraphics[width=\textwidth]{escudoInfor}
  \end{minipage}
  \vfill
  % Datos de alumno, curso y tutores ------------------
  \begin{center}%
  {%
    \noindent\LARGE
    Presentado por \@author{}\\ 
    en Universidad de Burgos --- \@date{}\\
    Tutor: \@tutor{}\\
  }%
  \end{center}%
  \null
  \cleardoublepage
  }
\makeatother


% Datos de portada
\title{Herramienta para la evaluación automática de entregas de programación con rúbricas y LLMs: IAGScore \\Documentación Técnica}
\author{Pedro Antonio Abellaneda Canales}
\tutor{Raúl Marticorena Sánchez}
\date{\today}

\begin{document}

\maketitle



\cleardoublepage



%%%%%%%%%%%%%%%%%%%%%%%%%%%%%%%%%%%%%%%%%%%%%%%%%%%%%%%%%%%%%%%%%%%%%%%%%%%%%%%%%%%%%%%%



\frontmatter


\clearpage

% Indices
\tableofcontents

\clearpage

\listoffigures

\clearpage

\listoftables

\clearpage

\mainmatter

\appendix

\apendice{Plan de Proyecto Software}

\section{Introducción}

\section{Planificación temporal}

\section{Estudio de viabilidad}

\subsection{Viabilidad económica}

\subsection{Viabilidad legal}



\apendice{Especificación de Requisitos}

\section{Introducción}

Una muestra de cómo podría ser una tabla de casos de uso:

% Caso de Uso 1 -> Consultar Experimentos.
\begin{table}[p]
	\centering
	\begin{tabularx}{\linewidth}{ p{0.21\columnwidth} p{0.71\columnwidth} }
		\toprule
		\textbf{CU-1}    & \textbf{Ejemplo de caso de uso}\\
		\toprule
		\textbf{Versión}              & 1.0    \\
		\textbf{Autor}                & Alumno \\
		\textbf{Requisitos asociados} & RF-xx, RF-xx \\
		\textbf{Descripción}          & La descripción del CU \\
		\textbf{Precondición}         & Precondiciones (podría haber más de una) \\
		\textbf{Acciones}             &
		\begin{enumerate}
			\def\labelenumi{\arabic{enumi}.}
			\tightlist
			\item Pasos del CU
			\item Pasos del CU (añadir tantos como sean necesarios)
		\end{enumerate}\\
		\textbf{Postcondición}        & Postcondiciones (podría haber más de una) \\
		\textbf{Excepciones}          & Excepciones \\
		\textbf{Importancia}          & Alta o Media o Baja... \\
		\bottomrule
	\end{tabularx}
	\caption{CU-1 Nombre del caso de uso.}
\end{table}

\section{Objetivos generales}

\section{Catálogo de requisitos}

\section{Especificación de requisitos}



\apendice{Especificación de diseño}

\section{Introducción}

\section{Diseño de datos}

\imagen{diagrama_ER.png}{Diagrama ER}
\imagen{diagrama_ER_2.png}{Diagrama de relación}



\section{Diseño arquitectónico}

\section{Diseño procedimental}




\apendice{Documentación técnica de programación}

\section{Introducción}

\section{Estructura de directorios}

\section{Manual del programador}

\section{Compilación, instalación y ejecución del proyecto}

\section{Pruebas del sistema}

\apendice{Documentación de usuario}

\section{Introducción}

En este apartado se ofrece una guía básica para que el usuario pueda interactuar de forma adecuada con la herramienta. Se describen los formularios principales y las distintas opciones disponibles para el usuario final al utilizar este entorno.

\section{Requisitos de usuarios}

Este proyecto ha sido diseñado para ejecutarse en entornos locales, con el objetivo de proteger la privacidad y seguridad de los datos personales generados a partir de las tareas del alumnado. Debido a la naturaleza de los servicios que integra (\textit{Ollama, Celery, Redis, PostgreSQL y Django}) no ha sido posible desplegarlo en plataformas gratuitas como \textit{Render o Heroku}, ya que presentan incompatibilidades técnicas. No obstante, el usuario puede instalar y ejecutar la aplicación en su propio equipo, siempre que cuente con \textit{Docker y Docker Compose} previamente instalados.

\section{Instalación}

Para realizar la instalación del proyecto, basta con seguir las instrucciones de detalladas en la 
sección \hyperref[sec:instalacion]{Compilación, instalación y ejecución del proyecto}


\section{Manual del usuario}

En esta sección se documenta el funcionamiento y la interacción con las principales características de la aplicación.

\subsubsection{Recursos de ejemplo}
El repositorio incluye una carpeta denominada \texttt{ejemplos}, que contiene recursos esenciales para facilitar la interacción con la herramienta. En esta carpeta se encuentran archivos que permiten:
Rellenar automáticamente un \textit{prompt} de ejemplo.
Importar una rúbrica predefinida.
Cargar tareas de muestra para realizar pruebas funcionales.
Estos recursos están diseñados para ayudar al usuario a familiarizarse rápidamente con el entorno y comprender su funcionamiento sin necesidad de generar contenido desde cero.

\subsubsection{Acceso a la aplicación}

Para acceder a la aplicación desde el navegador, simplemente introduce en la barra de direcciones la siguiente URL:

\begin{itemize}
    \item \href{http://localhost:8000}{\texttt{http://localhost:8000}}
\end{itemize}

Al ingresar, se mostrará la ventana principal desde la cual podrás iniciar sesión (\textit{login}). Para ello, deberás proporcionar tu \textit{email} y contraseña.

Si has creado un \textit{superusuario} siguiendo las indicaciones en la sección \hyperref[sec:superuser]{Crear \textit{superusuario}}, podrás acceder con esas credenciales y obtener acceso al panel de administración de \textit{Django}.

En caso de no disponer de credenciales, será necesario registrarse en la aplicación. Para ello, haz clic en la opción \textbf{Regístrate aquí} como se muestra en la Fig.~\ref{fig:registro_1}.

\imagen{registro_1}{Captura del botón de registro de usuario}

Se abrirá la ventana de registro Fig.~\ref{fig:view_registro}, donde deberás introducir un nombre, una dirección de \textit{email} y una contraseña. Además, deberás aceptar los Términos y condiciones y pulsar el botón \textbf{Crear una cuenta}.

\imagen{view_registro}{Captura de la ventana de registro de usuario}

Si el registro se ha completado correctamente, aparecerá un mensaje de confirmación Fig.~\ref{fig:registro_correcto}, como el siguiente:

\imagen{registro_correcto}{Captura del mensaje de registro correcto}

A continuación, podrás acceder a la página de bienvenida de la aplicación, Fig.~\ref{fig:home}. En esta página encontrarás una barra de navegación con las distintas secciones disponibles, botones de acceso rápido, un selector de idioma y un botón para cerrar sesión.

\imagen{home}{Captura de la página de bienvenida}

Desde esta página se puede seguir el flujo lógico propuesto a través de los menús, comenzando por la creación de un \textit{prompt}, la importación de una rúbrica, la configuración de una corrección mediante la importación de tareas y, finalmente, la ejecución del proceso de corrección.

Para realizar el primer paso, accede a la sección \textbf{Prompts}, bien a través del menú de navegación o mediante el botón de acceso rápido. En esta sección se mostrarán los \textit{prompts} ya creados (en caso de que existan) y se ofrecerá la opción de crear uno nuevo. Para ello, pulsa el botón \textbf{Nuevo Prompt}, Fig.~\ref{fig:nuevo_prompt}, introduce un nombre identificativo y el contenido del \textit{prompt}, donde se debe indicar al modelo que se le proporcionará una rúbrica y un conjunto de tareas para realizar la corrección.

\imagen{nuevo_prompt}{Captura de la creación de un nuevo \textit{prompt}}

Tras pulsar el botón \textbf{Guardar}, si todo ha funcionado correctamente, se mostrará un mensaje de confirmación indicando que el \textit{prompt} se ha creado con éxito y aparecerá listado en la tabla correspondiente, junto con la fecha de creación y las acciones disponibles, como se puede ver en Fig.~\ref{fig:prompt_creado}.

\imagen{prompt_creado}{Captura del \textit{prompt} creado correctamente}

El siguiente paso consiste en importar una rúbrica en formato Markdown. Para ello, debemos acceder a la sección \textbf{Rúbricas}.

Una vez en esta sección, se debe pulsar el botón \textbf{Nueva rúbrica}, lo que abrirá un formulario, como el que se muestra en Fig.~\ref{fig:nueva_rubrica}, donde podremos introducir un nombre identificativo e importar un archivo en formato \texttt{.md} que contenga el contenido de la rúbrica.

\imagen{nueva_rubrica}{Captura del formulario de importación de una rúbrica}

Al pulsar el botón \textbf{Guardar}, si el proceso se realiza correctamente, se mostrará un mensaje de confirmación, Fig.~\ref{fig:rubrica_creada}, indicando que la rúbrica ha sido creada con éxito. Esta aparecerá listada en la tabla correspondiente, junto con la fecha de creación y las acciones disponibles.

\imagen{rubrica_creada}{Captura de la rúbrica creada correctamente}

Importar la rúbrica en formato \textit{Markdown} resulta especialmente útil, ya que permite una mejor visualización del contenido estructurado, Fig.~\ref{fig:rubrica_html}. Esta visualización puede comprobarse pulsando el botón \textbf{Mostrar}.

\imagen{rubrica_html}{Captura mostrando rúbrica formateada en \textit{HTML}}

Una vez disponemos de un \textit{prompt} y de una rúbrica, ya podemos proceder a configurar una corrección. Para ello, accedemos a la sección \textbf{Correcciones}, seleccionando la opción correspondiente en la barra de navegación y, a continuación, haciendo clic en \textbf{Nueva corrección}, como se puede ver en Fig.~\ref{fig:new_correction}.

\imagen{new_correction}{Captura mostrando el botón \textbf{Nueva corrección}}

Se abrirá el formulario de configuración, Fig.~\ref{fig:config_correction}, donde será necesario seleccionar el \textit{prompt} y la rúbrica previamente creados. Además, se deberá cargar un archivo comprimido (en formato \texttt{.zip}, \texttt{.rar}, \texttt{.7z} o \texttt{.tar}) que contenga las tareas a corregir. También será preciso introducir un nombre o una breve descripción que identifique la corrección y escoger un modelo disponible en la lista (el cual debe haber sido previamente \hyperref[sec:pull_model]{cargado}).

A continuación, se mostrarán una serie de parámetros avanzados, tales como \textbf{Temperatura}, \textbf{Top\_p}, \textbf{Top\_k}, \textbf{Ventana de contexto} y \textbf{Formato de salida}, que vienen configurados por defecto, aunque el usuario podrá modificarlos según sus necesidades.

\imagen{config_correction}{Captura mostrando el formulario de configuración de corrección}

Una vez completados los datos, al pulsar el botón \textbf{Guardar configuración} la aplicación redirigirá automáticamente a la página principal de correcciones. Allí se listarán todas las correcciones configuradas, junto con información relevante y las acciones disponibles para cada una, como se muestra en Fig.~\ref{fig:create_correction}.

\imagen{create_correction}{Captura mostrando una corrección configurada}

Con la corrección configurada, estará lista para ejecutarse pulsando la opción \textbf{Ejecutar}. Al hacerlo, se mostrará un mensaje indicando que la corrección está en proceso y que el resultado podrá consultarse más adelante. Esta operación se realiza en segundo plano, por lo que no bloquea la aplicación. Mientras la corrección se ejecuta, su estado aparecerá como \textit{Ejecutando...} y se deshabilitarán las acciones sobre ella para garantizar la consistencia del sistema, como puede verse en Fig.~\ref{fig:run_correction}.

\imagen{run_correction}{Captura mostrando una corrección en ejecución}

Al finalizar el proceso, estará disponible la opción para descargar la corrección, Fig.~\ref{fig:descarga}, en formato \texttt{.txt}.

\imagen{descarga}{Captura mostrando la opción de descarga}

De este modo, se completa un ciclo completo de la funcionalidad principal de la aplicación.

La sección \textbf{\textit{LLM}} ofrece detalles sobre el modelo utilizado en las diversas pruebas realizadas, mientras que la sección \textbf{Documentación} contiene la documentación técnica del proyecto, junto con el manual de instalación.

\apendice{Anexo de sostenibilización curricular}

\section*{F.1. Introducción}
Este anexo expone cómo el desarrollo de la herramienta \textit{IAGScore}, presentada en este Trabajo de Fin de Grado, se alinea con los Objetivos de Desarrollo Sostenible (\textit{ODS}) propuestos por la Agenda 2030 de las Naciones Unidas~\cite{web:agenda2030}. La automatización de la evaluación de tareas de programación mediante modelos de lenguaje (\textit{LLMs}) no solo supone un avance técnico, sino que también permite avanzar hacia una educación más justa, eficiente, privada y accesible.

\section*{F.2. Objetivos de Desarrollo Sostenible relacionados}

\textbf{\textit{ODS} 4 – Educación de calidad:}
Este objetivo~\cite{wiki:ods4} busca garantizar una educación inclusiva, equitativa y de calidad, promoviendo el aprendizaje a lo largo de la vida. \textit{IAGScore} contribuye a este fin al ofrecer una evaluación más justa y objetiva, reduciendo los sesgos de la corrección manual. Al funcionar localmente, protege los datos personales del alumnado, asegurando la confidencialidad y el cumplimiento de la normativa de protección de datos. Esto permite un uso ético y respetuoso de los \textit{LLMs} en el ámbito educativo.

\textbf{\textit{ODS} 10 – Reducción de las desigualdades:~\cite{wiki:ods10}}
La evaluación automatizada con criterios unificados y trazables promueve la equidad educativa, garantizando una valoración coherente del esfuerzo estudiantil, sin importar el contexto institucional. Al no depender de servicios en la nube, \textit{IAGScore} evita sesgos externos y barreras tecnológicas, facilitando su adopción en centros con distintos niveles de recursos.

\textbf{\textit{ODS} 9 – Industria, innovación e infraestructura:}  
este \textit{ODS}~\cite{wiki:ods9} promueve la construcción de infraestructuras resilientes, la industrialización sostenible y la innovación. \textit{IAGScore} representa una innovación aplicada al sector educativo al integrar tecnologías emergentes como los \textit{LLMs} en procesos institucionales clave. La posibilidad de desplegar estos modelos en entornos locales es una ventaja estratégica, ya que permite mantener el control completo sobre la infraestructura, reducir la dependencia de servicios externos y evitar la externalización de datos sensibles. Además, el uso de tecnologías abiertas (licencias \textit{MIT} y \textit{BSD}) fomenta la creación de una infraestructura tecnológica accesible, reproducible y escalable.

\section*{F.3. Accesibilidad y adherencia tecnológica}
Desde el diseño inicial del proyecto se ha priorizado la usabilidad de la herramienta, con una interfaz sencilla que permite al profesorado sin conocimientos técnicos integrar rúbricas personalizadas y evaluar tareas mediante simples interacciones. La facilidad de uso, la documentación técnica incluida y la transparencia del proceso de corrección aumentan la adopción y fidelización del usuario (profesorado), lo cual es fundamental para que su aplicación tenga impacto real.

La posibilidad de ejecutar la herramienta en máquinas locales o servidores institucionales garantiza una accesibilidad adecuada, incluso en entornos con limitaciones de conectividad o políticas estrictas de protección de datos. Además, se evitan riesgos asociados a la transferencia de datos personales fuera del entorno académico, lo cual refuerza la confianza en su adopción y cumplimiento normativo.

\section*{F.4. Futuras líneas de trabajo}
Aunque el enfoque actual de la herramienta se centra en mejorar la eficiencia y objetividad de la evaluación, se prevé incorporar mecanismos explícitos que informen a los usuarios sobre el impacto de su uso en términos de sostenibilidad. Por ejemplo:

\begin{itemize}
  \item Incorporar estadísticas de uso y ahorro de tiempo, evidenciando el impacto positivo en la gestión docente.
  \item Integrar análisis de sesgo para garantizar que los modelos empleados mantengan criterios justos y no discriminatorios.
  \item Ampliar las opciones de despliegue local con soporte para arquitecturas ligeras y entornos desconectados de internet.
  \item Implementar una API REST que permita integrar la herramienta con plataformas educativas existentes como Moodle, lo que facilitaría su adopción en contextos institucionales amplios, automatizando aún más los flujos de evaluación.
\end{itemize}


\section*{F.5. Conclusión}

Este TFG demuestra que es posible alinear el desarrollo tecnológico con los principios de sostenibilidad. El diseño y aplicación de \textit{IAGScore} no solo ha implicado retos técnicos, sino también decisiones conscientes orientadas a promover una educación más equitativa, eficiente, privada y responsable. En un contexto donde el acceso a la tecnología puede profundizar o mitigar desigualdades, herramientas como esta representan un paso hacia una transformación digital verdaderamente inclusiva y respetuosa con los derechos del alumnado.

No obstante, es importante reconocer que el uso de modelos de lenguaje de gran escala (LLMs) conlleva un elevado consumo energético, lo que plantea desafíos en términos de sostenibilidad ambiental. Aunque \textit{IAGScore} cumple con varios Objetivos de Desarrollo Sostenible, este aspecto debe ser considerado en futuras versiones, buscando un equilibrio entre innovación educativa y responsabilidad ecológica.


\bibliographystyle{plain}
\bibliography{bibliografiaAnexos}

\end{document}
